\clearpage
\section{Historischer Abriss}
\label{cuts:history}

Richard Dedekinds Schrift \emph{Stetigkeit und irrationale Zahlen}
erschien 1872. Im Vorwort führt er die Überlegungen auf seine Zürcher
Vorlesungstätigkeit im Jahr 1858 zurück: Die Analysis sollte eine
arithmetische Begründung erhalten, die sich nicht auf eine als lückenlos
vorgestellte geometrische Gerade stützt.

In §~4 beschreibt Dedekind einen Schnitt durch zwei Klassen rationaler
Zahlen: Jede Zahl der unteren Klasse liegt unter jeder Zahl der oberen.
Schnitte, die keine rationale Zahl hervorbringt, führen ihn zu neuen
irrationalen Zahlen. In §~5 zeigt er die Stetigkeit des so erweiterten
Zahlbereichs. Seine Behandlung der Rechenoperationen folgt erst danach.

Unsere Darstellung verwendet eine heute übliche Festlegung: Nur die
untere Klasse wird gespeichert, und sie besitzt kein größtes Element.
Dadurch entfällt die doppelte Darstellung eines rationalen Randes.
Die Vollständigkeit formulieren wir als Supremumseigenschaft;
die Vereinigung liefert ihre konkrete Konstruktion.

\subsection*{Originalquelle}
\begingroup\small
Richard Dedekind: \emph{Stetigkeit und irrationale Zahlen}.
Braunschweig: Friedrich Vieweg und Sohn, 1872.
Hier besonders Vorwort und §§~4--5.\par
\href{https://doi.org/10.3931/e-rara-71757}{Digitalisat der Erstausgabe, ETH-Bibliothek Zürich, Rar 22802}.\par
\href{https://www.math.ru.nl/werkgroepen/gmfw/bronnen/dedekind2.html}{Transkription der vierten Auflage von 1912, Radboud Universiteit}.
\endgroup

\clearpage
\section{Schlussbemerkung}
\label{cuts:conclusion}

Der Weg begann mit einer einfachen Beschreibung: Eine Zahl wird durch
die rationalen Zahlen unter ihr bestimmt. Vier Bedingungen machen
aus dieser Vorstellung einen Schnitt. Die Inklusion liefert seine Ordnung;
die rationalen Hauptschnitte bewahren die vertraute rationale Ordnung.

Für die Vollständigkeit brauchten wir dann nur einen neuen Kandidaten:
die Vereinigung einer nichtleeren, nach oben beschränkten Schnittfamilie.
Sie enthält alles Nötige und nichts, was eine obere Schranke zusätzlich
fordern müsste. Sobald ihre Schnittbedingungen feststehen, ist der
Supremumsbeweis nur noch die Erklärung dieses Satzes.

\begin{center}
\fbox{\begin{minipage}{.86\linewidth}
\centering
\textbf{Zahlen als untere Mengen.}\par\smallskip
Vergleichen heißt einschließen.\par
Ein Supremum bilden heißt vereinigen.
\end{minipage}}
\end{center}

Damit ist die vollständige Ordnung konstruiert. Addition und
Multiplikation der Schnitte können daran anschließend entwickelt werden;
sie gehören zum weiteren Aufbau von Band~19.
Hier reicht bereits die Ordnung für den entscheidenden Gewinn:
Nichtleere beschränkte Familien haben eine kleinste obere Schranke,
und wir wissen ausdrücklich, wie sie entsteht.

Das Übungsheft führt noch einmal durch die tragenden Schritte.
Besonders aufschlussreich ist es, in jedem Beweis zu markieren,
welche der vier Schnittbedingungen tatsächlich benutzt wird.
